\documentclass[a4paper]{article}

\usepackage{amsthm,amsmath,mathtools,gensymb,braket}

\title{July 14, 2026}
\author{Miel Gabriel Lagdan}

\begin{document}
	\maketitle
	Here, we introduce a potential substitution to the differential quotient of the first type.
	\section{The Differential Quotient of First Type}
	Born and Jordan on their paper "On Quantum Mechanics" provided an algorithm in finding the derivative of an operator function, that is
	\begin{equation}\label{eqn:dqft}
		\left(\frac{\partial\textbf{y}}{\partial\textbf{y}_k}\right)_1=\sum_{r=1}^s\delta_{l_r,k}\prod_{m=r+1}^s\textbf{y}_{l_m}\prod_{m=1}^{r-1}\textbf{y}_{l_m}
	\end{equation}
	that is the algorithm dictates that
	\begin{enumerate}
		\item For a given $r$, it verifies if $\textbf{y}_k=\textbf{y}_{l_r}$, if so, this makes $\delta_{l_r,k}=1$.
		\item Every operator element after $r$ is copied. Then every operator element before $r$ is the copied after.
	\end{enumerate}
	
	Bagunu and Galapon\footnote{Bagunu \& Galapon. Quantization of the Hamilton Equations of Motion.} was able to verify that any of the quantization images (that is \textit{Weyl}, \textit{simplest symmetrization}, and \textit{Born-Jordan}) will result in Born-Jordan images because of its basis. The paper was also able to modify Eq. \ref{eqn:dqft} to result in Weyl and simplest-symmetrization images.
	
	The paper also made a claim that due to the case-by-case modification of the differential quotient of first type, there is no method to generalize the differential quotient of first type.
	
	This investigation attempts to provide another formulation for the differential quotient of first type, by for the purposes of this investigation will be referred to as the DQST.
	
	\section{Attempts to Differentiate Operators}
	Neudecker\footnote{Neudecker. Some Theorems on Matrix Differentiation with Special Reference to Kronecker Matrix Products} writes the following theorems must be considered to work with matrix differentials, by which we consider our operators to be matrices.
	\begin{itemize}
		\item They obey \textbf{Leibniz Rule}. That is, for conformable matrices $X$ and $Y$, $d(XY)=(dX)Y+X(dY)$
		\item If $l$ is a linear scalar function of the matrix $X$, then $d(lX)=l(dX)$
		\item $d~\text{vec}~X=\text{vec}~(dX)$
	\end{itemize}
\end{document}